
    \documentclass[10pt,a4paper,twocolumn]{article}
    \usepackage[margin=1.8cm]{geometry}
    \usepackage{booktabs}
    \usepackage{graphicx}
    \usepackage{xcolor}
    \usepackage{hyperref}
    \usepackage[small]{titlesec}
    \setlength{\parindent}{0pt}
    \setlength{\parskip}{0.4em}

    \title{\textbf{Evidence Run Report}\\[0.3em]
           \large 2026-03-12\_evidence-v1.5}
    \author{autoresearcher2}
    \date{2026-03-12 22:49}

    \begin{document}
    \maketitle
    \thispagestyle{empty}

    \section*{Status: IN PROGRESS (autoresearch, full running)}

    Head-to-head comparison of four experiment-selection approaches on real \texttt{train.py}.
    Schema: DEPTH $\times$ MATRIX\_LR $\times$ WEIGHT\_DECAY (27 cells).
    Budget: 20 experiments per approach. Hardware: GPU~1 (RTX PRO 6000), seed~42.

    \section{Results}

    \begin{table}[h]
    \centering
    \small
    \begin{tabular}{lrrrrr}
    \toprule
    Approach & Done & Best & Mean & Cells & Tok/exp \\
    \midrule
        random & 20/20 & 1.033551 & 1.054828 & 15 & 179 \\
    bayesian & 20/20 & 1.033778 & 1.051689 & 20 & 205 \\
    autoresearch & 19/20 & 1.032934 & 1.046333 & 15 & 183 \\
    full & 13/20 & 1.033423 & 1.048775 & 13 & 204 \\
    \bottomrule
    \end{tabular}
    \caption{Performance summary. Best/mean = val\_bpb (lower is better). Cells = unique cells explored. Tok/exp = mean tokens trained per experiment (millions).}
    \end{table}

    \section{Convergence}

    Experiment number at which each approach first reached within 0.001 of its final best:

    \begin{table}[h]
    \centering
    \small
    \begin{tabular}{lrr}
    \toprule
    Approach & Conv.\ exp & Final best \\
    \midrule
        random & 11 & 1.033551 \\
    bayesian & 5 & 1.033778 \\
    autoresearch & 8 & 1.032934 \\
    full & 3 & 1.033423 \\
    \bottomrule
    \end{tabular}
    \end{table}

    \section{Decision Source Integrity}

    Every experiment carries a source label indicating what method actually chose it.
    \textcolor{red}{\textbf{WARNING: 5 fallback(s) detected.}}

    \begin{table}[h]
    \centering
    \small
    \begin{tabular}{ll}
    \toprule
    Approach & Sources \\
    \midrule
        random & random: 20 \\
    bayesian & lookahead: 20 \\
    autoresearch & llm\_fallback\_empty: 3, llm\_flat: 16, random\_init: 1 \\
    full & llm\_augmented: 11, llm\_fallback: 2, lookahead: 7 \\
    \bottomrule
    \end{tabular}
    \end{table}

    \section{Observations (not claims)}

    These are patterns visible in the data. They are \textbf{not} generalization claims ---
    this is one run with one seed on one schema.

    \begin{itemize}
    \item DEPTH=8 consistently produces the best val\_bpb across all approaches.
          DEPTH=10 is worst (fewer tokens trained in fixed wall time).
    \item WEIGHT\_DECAY=0.4 appears in most best-performing configurations.
    \item Reproducibility is strong: repeated cell visits show $\Delta$val\_bpb $< 0.001$.
    \item Throughput varies $\sim$3$\times$ across depths (300K--1M tok/s), meaning
          deeper models train fewer tokens in the same wall time.
    \end{itemize}

    \section{What this does NOT prove}

    \begin{itemize}
    \item Generalization beyond this schema, seed, or hardware.
    \item That any approach is \emph{better} in a statistically significant sense
          (single seed, small budget).
    \item That the LLM's advantage (if confirmed) comes from world knowledge vs.\ lucky guessing.
    \item That structured signals (appraisal, factor importances) help the LLM.
    \end{itemize}

    \section{Methodology notes}

    Random and bayesian ran in the original process (run\_id 5d2ffa21). A silent-fallback
bug was discovered: LLM failures were labeled as "llm\_flat" instead of "llm\_fallback\_*".
Process was killed after bayesian completed. Autoresearch and full were restarted
with fixed source-labeling code (run\_id 65018a7c, commit ce7d19a).


    \section{Discussion}
    The headline result is that all four approaches converge to similar best val\_bpb values (1.0329--1.0338), a spread of less than 0.001. In a 27-cell grid with one dominant factor, this is unsurprising: even uniform random search covers 15 of 27 cells in 20 trials (observed), giving it a $>$99\% chance of hitting the best region. \textbf{The search space is too small to discriminate the approaches on best-found performance.}

Mean val\_bpb tells a more interesting story. Ranking by mean: autoresearch (1.046) $<$ full (1.049) $<$ bayesian (1.052) $<$ random (1.055). The LLM-containing approaches waste fewer experiments on poor regions (\emph{observed}). This is consistent with the hypothesis that LLM priors help allocate trials, but the effect is modest---roughly 0.006 bpb between autoresearch and bayesian---and drawn from a single seed with no significance testing.

\textbf{The Bayesian model did not clearly help.} Bayesian alone (mean 1.052) barely beats random (1.055), and full (Bayesian + LLM, mean 1.049) did not beat autoresearch (LLM alone, mean 1.046). If anything, the Bayesian component adds overhead without improving selection quality. Two candidate explanations: (i) a linear factor model over one-hot features is too crude to capture the depth$\times$learning-rate interaction that dominates this space (\emph{plausible}); (ii) with only 27 cells, Thompson sampling has little room to outperform random before the space is saturated (\emph{supported} by the coverage numbers---bayesian hit 20/27 cells, suggesting it explored broadly rather than exploiting).

\textbf{Source integrity raises a concern for the full approach.} Only 13 of 20 experiments completed, and sources are split: 7 lookahead, 11 llm\_augmented, 2 llm\_fallback. The full approach is not purely ``Bayesian + LLM''---it is a mixture that fell back to the Bayesian arm roughly a third of the time. Similarly, autoresearch had 3/19 fallbacks. These fallbacks are honestly labeled, but they weaken the factorial interpretation: the 2$\times$2 design is confounded by variable fallback rates.

\textbf{The token budget confound is unresolved.} DEPTH=6 trains $\sim$3$\times$ more tokens than DEPTH=10 in the same wall time. All approaches converged on DEPTH=8 as optimal, which trains $\sim$1.8$\times$ fewer tokens than DEPTH=6. This suggests DEPTH=8 is genuinely better architecture, not merely better-trained (\emph{supported}), but we cannot fully separate architecture quality from training duration without controlled token budgets.

\begin{itemize}
\item \textbf{What we learned:} LLM priors improve \emph{average} trial quality in this small grid. The Bayesian model adds no detectable value at this scale. The combination does not outperform the LLM alone.
\item \textbf{What we don't know:} Whether the Bayesian component helps in larger search spaces where LLM priors degrade and systematic uncertainty tracking matters. Whether autoresearch's advantage holds across seeds.
\item \textbf{Next steps:} (1) Repeat with a larger schema ($\geq$100 cells) where random search cannot saturate the space. (2) Run multiple seeds to get confidence intervals. (3) Controlled token-budget experiments (fixed tokens, not fixed wall time). (4) Complete the full approach's remaining 7 experiments before drawing final conclusions.
\end{itemize}

    \begin{figure*}[b]
    \centering
    \includegraphics[width=\textwidth]{/Users/erik/github.com/erikdebruijn/autoresearcher2/artifacts/runs/2026-03-12_evidence-v1.5/figures/evidence_dashboard.png}
    \caption{Evidence dashboard. Top-left: convergence curves. Top-right: factor effects.
    Bottom: exploration heatmaps per approach (cell visit counts).}
    \end{figure*}
    
    \end{document}
    